% Options for packages loaded elsewhere
\PassOptionsToPackage{unicode}{hyperref}
\PassOptionsToPackage{hyphens}{url}
\PassOptionsToPackage{dvipsnames,svgnames,x11names}{xcolor}
\documentclass[
  10pt,
  fontset=fandol]{ctexart}
\usepackage{xcolor}
\usepackage[margin=16mm]{geometry}
\usepackage{amsmath,amssymb}
\setcounter{secnumdepth}{-\maxdimen} % remove section numbering
\usepackage{iftex}
\ifPDFTeX
  \usepackage[T1]{fontenc}
  \usepackage[utf8]{inputenc}
  \usepackage{textcomp} % provide euro and other symbols
\else % if luatex or xetex
  \usepackage{unicode-math} % this also loads fontspec
  \defaultfontfeatures{Scale=MatchLowercase}
  \defaultfontfeatures[\rmfamily]{Ligatures=TeX,Scale=1}
\fi
\usepackage{lmodern}
\ifPDFTeX\else
  % xetex/luatex font selection
\fi
% Use upquote if available, for straight quotes in verbatim environments
\IfFileExists{upquote.sty}{\usepackage{upquote}}{}
\IfFileExists{microtype.sty}{% use microtype if available
  \usepackage[]{microtype}
  \UseMicrotypeSet[protrusion]{basicmath} % disable protrusion for tt fonts
}{}
\makeatletter
\@ifundefined{KOMAClassName}{% if non-KOMA class
  \IfFileExists{parskip.sty}{%
    \usepackage{parskip}
  }{% else
    \setlength{\parindent}{0pt}
    \setlength{\parskip}{6pt plus 2pt minus 1pt}}
}{% if KOMA class
  \KOMAoptions{parskip=half}}
\makeatother
\setlength{\emergencystretch}{3em} % prevent overfull lines
\providecommand{\tightlist}{%
  \setlength{\itemsep}{0pt}\setlength{\parskip}{0pt}}
\usepackage{amsmath,amssymb,mathtools,needspace}
\xeCJKDeclareCharClass{CJK}{"2032,"0400->"04FF,"2160->"216B,"2460->"2473,"25A0,"25C7,"25CB,"25CF}
\IfFontExistsTF{Arial Unicode MS}{\xeCJKsetup{AutoFallBack=true}\setCJKfallbackfamilyfont{\CJKrmdefault}{Arial Unicode MS}}{}
\setcounter{secnumdepth}{0}
\setlength{\emergencystretch}{2em}
\allowdisplaybreaks
\usepackage{bookmark}
\IfFileExists{xurl.sty}{\usepackage{xurl}}{} % add URL line breaks if available
\urlstyle{same}
\hypersetup{
  pdftitle={数列求和的累加算法},
  colorlinks=true,
  linkcolor={Maroon},
  filecolor={Maroon},
  citecolor={Blue},
  urlcolor={Blue},
  pdfcreator={LaTeX via pandoc}}

\title{数列求和的累加算法}
\author{}
\date{}

\begin{document}
\maketitle

\subsubsection{13.2.1 数列求和的累加算法}\label{ux6570ux5217ux6c42ux548cux7684ux7d2fux52a0ux7b97ux6cd5}

下述数列求和问题是人们所熟知的：

\[
S=a_0+a_1+\cdots+a_n.
\tag{13.2.1}
\]

这个计算模型有两个简单的特例。当 \(n=0\) 即为一项和式 \(S=a_0\) 时，所给计算模型就是它的解，这时不需要做任何计算。这表明，对于数列求和问题，它的解是计算模型退化的情形。又当 \(n=1\) 即计算两项和式 \(S=a_0+a_1\) 时，计算过程是平凡的，这时不存在算法设计问题。

现在基于这两种简单情形考察所给和式（13.2.1）的累加求和算法。设 \(b_k\) 表示前 \(k+1\) 项的部分和 \(a_0+a_1+\cdots+a_k\)，则有

\href{../../source-images/pdf-309.jpeg}{原页图像}

\[
\begin{cases}
b_0=a_0,\\
b_k=b_{k-1}+a_k,\quad k=1,2,\cdots,n,
\end{cases}
\tag{13.2.2}
\]

而计算结果 \(b_n\) 即为所求的和值

\[
S=b_n.
\tag{13.2.3}
\]

上述\textbf{数列求和的累加算法，其设计思想是将多项求和（式（13.2.1））化归为两项求和（式（13.2.2））的重复。}而依式（13.2.2）重复加工若干次，最终即可将所给和式（13.2.1）加工成一项和式（13.2.3）的退化情形，从而得出和值 \(S\)。

再剖析计算模型自身的演变过程。按式（13.2.2）每加工一次，所给和式（13.2.1）便减少一项，而所生成的计算模型依然是数列求和。反复施行这种加工手续，计算模型不断变形为

\[
\underset{\text{（计算模型）}}{n+1\text{ 项和式}}
\Rightarrow n\text{ 项和式}
\Rightarrow n-1\text{ 项和式}
\Rightarrow\cdots\Rightarrow
\underset{\text{（所求结果）}}{1\text{ 项和式}},
\]

这里，符号``\(\Rightarrow\)''表示重复施行两项求和的加工手续。

这样，如果定义和式的项数为数列求和问题的\textbf{规模}，则所求和值可以视作规模为 1 的退化情形。因此，只要令和式的规模（项数）逐次减 1，最终当规模为 1 时即可直接得出所求的和值。这样设计出的算法就是累加求和算法。

上述累加求和算法可以视作规模缩减技术的一个范例。

\end{document}
